\chapter*{Abstract}
\addcontentsline{toc}{chapter}{Abstract}
Bangladesh’s rapid, largely unregulated industrialization has resulted in major air pollution hotspots in factories and workshops, including welding bays, boiler rooms, generator rooms, chemical storage areas and factory entrances. Workers may be exposed to high concentrations of particulate matter (PM1.0, PM2.5, PM10), carbon monoxide (CO), nitrogen dioxide (NO$_2$), volatile organic compounds (VOC), ammonia (NH$_3$), hydrogen sulfide (H$_2$S) and smoke. Conventional air quality management is mainly on the city scale or plant scale, and centralized purification or coarse monitoring cannot effectively protect workers near specific pollution sources. This capstone presents \textbf{NirmalNode}, a Green Internet of Things (IoT) and Adaptive Artificial Intelligence (AI) powered system for local, hotspot level air purification.

NirmalNode is built around a low-cost ESP32 microcontroller and can be equipped with a PMS5003 particulate sensor, MiCS-4514 and MQ-series gas sensors, and optionally a BME280 environmental sensor. The system uses Wi-Fi to send information about pollutants and the environment, including time and location, to a cloud database. A lightweight adaptive AI model that can be deployed at the edge to predict pollutant concentrations one to two hours ahead by leveraging historical Bangladesh air-quality data and incremental (online) learning. If the forecast exceeds a hazard threshold, a DC fan is turned on to pull polluted air through a cyclone separator, HEPA filter and activated-carbon filter, before returning purified air to the immediate breathing zone of the worker.

The main methodological contribution is the use of incremental/online learning models, such as SGD Regressor, Passive-Aggressive Regressor, Hoeffding Tree and Adaptive Random Forest, implemented in River and scikit-learn. These models enable NirmalNode to adapt to different industrial environments without the need for periodic full retraining and manual redeployment. The area-specific adaptivity combined with lightweight Green AI models and low-power ESP32-centred hardware, paves a way for cost-effective deployment in small and medium industrial units of Bangladesh.

A bench-scale prototype was designed and tested in a controlled laboratory environment and, where possible, at a representative industrial site. The prediction accuracy was assessed using MAE, RMSE, Precision, Recall, F1-score and $R^2$. The purification efficiency was assessed in terms of pollutant reduction and power consumption and estimated deployment cost. This study fills the research gap of a low-cost, adaptive, hotspot-targeted purification system for Bangladesh. Future work includes multi-node factory deployment, federated learning, further edge-AI optimization and a solar powered variant for areas with unreliable grid power.

\vspace{0.5cm}
\textbf{Keywords:} Green IoT, Green AI, Adaptive AI, Incremental Learning, Air Pollution Prediction, Local Air Purification, ESP32, HEPA Filter, Activated Carbon, Cyclone Separator, Industrial Hotspot Monitoring, Edge-Cloud Architecture, Bangladesh
